\section{中心三次対の全Artin形式局所構造と二次包絡}
\label{sec:moduli-artin}
本節は、滑らかな可換三次三様体の近傍において、
有限係数の17次元族と異なる\emph{完備形式局所モジュライ}を扱う。
生成空間の座標標識を残す関手と、その射影座標変更を割った関手を
区別する。構造代数層との比較に用いるのは後者であり、
特殊ファイバーの同一視は保持する。主張は通常の集合値Artin変形関手に関するものであり、
derived stackの同型でも非可換領域全体の大域分類でもない。
既知のKontsevich量子化から中心三次商の族を実際に構成し、
構造代数層の変形関手と比較する。
全ての次数付き代数変形を同次formalityで先に分類することは仮定しない。
構造代数層への移行ではYekutieliの量子化定理の仮定と
相対コホモロジー計算を明示する。

以下$k=\C$、$V=k^5$、$S=\Sym V^*$、
$F\in S_3$を滑らかな三次三様体
$X=V_+(F)\subset\PP(V)$の定義式とする。
$(R,\mathfrak m)$は剰余体$k$の局所Artin代数である。

\begin{lemma}[平坦な正規三次商は自動的に中心商]
\label{lem:artin-normal-central}
$B_R$を$R$上各次数有限自由な連結非負次数付き代数とし、
一次で生成され、特殊ファイバー$B_R/\mathfrak mB_R=S$とする。
$w\in(B_R)_3$が$F$に特殊化するなら$w$は左右正則であり、
以下の条件は同値である：
\[
 B_R/(w)\text{ が各次数で }R\text{-平坦},\qquad
 wB_R=B_Rw,\qquad
 w\in Z(B_R).
\]
成立するとき商のHilbert級数は$(1-t^3)/(1-t)^5$である。
\end{lemma}
\begin{proof}
$w$の左右乗法の各次数成分は、$\mathfrak m$で還元すると
$F:S_{d-3}\hookrightarrow S_d$である。
$F\ne0$と$S$が整域であることから単射となり、
有限自由行列の最大rank小行列式が単元であるため、
元の乗法も分裂単射になる。

$J=(w)$について商が平坦なら
$0\to J_d\to(B_R)_d\to(B_R/J)_d\to0$は
有限自由$R$加群の分裂完全列である。
$J_d/\mathfrak mJ_d=FS_{d-3}$なので、
左右乗法$(B_R)_{d-3}\to J_d$はNakayamaと
先の単射性により同型である。従って$J=wB_R=B_Rw$。
正規性から$wb=\sigma(b)w$で次数付き自己同型$\sigma$が得られる。
特殊ファイバーでは$\bar\sigma=\mathrm{id}$であり、
$b=w$と右消去によって$\sigma(w)=w$となる。
$\sigma\equiv\mathrm{id}\pmod{\mathfrak m^r}$とし、
$\mathfrak m^r/\mathfrak m^{r+1}=I$とおくと、
次の差分は次数零の導分
\[
 D:S\longrightarrow I\otimes_k S,\qquad
 D(F)=0
\]
である。$D=\sum_i\ell_i\partial_{x_i}$、
$\ell_i\in I\otimes S_1$だから
$\sum_i\ell_iF_i=0$。滑らかさから
$F_0,\ldots,F_4$は二次正則列をなし、
Koszul複体に一次係数のsyzygyが存在しない。
よって$D=0$となり、$\mathfrak m$の冪零性から
$\sigma=\mathrm{id}$を得る。
中心性から商平坦性へは、
各次数の分裂単射$w:(B_R)_{d-3}\to(B_R)_d$を使えばよい。
\end{proof}

\begin{corollary}[標識付き対の自己同型]
\label{cor:artin-scalar}
上の中心正規対$(B_R,Rw)$の次数付き自己同型が
特殊ファイバーで恒等なら、次数一の生成元への作用は
すべて同じ単元$\lambda\in1+\mathfrak m$によるスカラー倍である。
従って射影標識を入れた対には非自明な無限小自己同型はない。
\end{corollary}
\begin{proof}
$\phi(w)=cw$、$c\in1+\mathfrak m$と書き、
$c$の一意の三乗根による次数スカラー倍を合成すると
$w$を固定する自己同型に帰着する。
その$\mathfrak m$進差分は前補題の$D(F)=0$と同じ
一次Koszul syzygyで消える。
\end{proof}

\subsection{相対Casimir条件と五つのPl\"ucker方程式}
体積形式$\Omega=dx_0\wedge\cdots\wedge dx_4$を固定する。
$F_R\in R\otimes S_3$を$F$の持上げとし、
$a=\sum_{i<j}a_{ij}\,dx_i\wedge dx_j$を
係数$a_{ij}\in\mathfrak m$の定数二形式とする。
\[
 \iota_{\Pi(F_R,a)}\Omega=dF_R\wedge a
\]
で二次bivectorを定める。
$\{i,j,k,l\}=\{0,\ldots,4\}\setminus\{r\}$、
$i<j<k<l$に対し
\[
 q_r(a)=(-1)^r
 (a_{ij}a_{kl}-a_{ik}a_{jl}+a_{il}a_{jk}),\qquad
 g=\sum_{r=0}^4q_r(a)\,\partial_rF_R
\]
とおく。

\begin{lemma}[全Artin上でのCasimir核]
\label{lem:artin-casimir}
二次bivector$\Pi$が$\iota_{dF_R}\Pi=0$を満たすことと、
一意な$a\in R\otimes\wedge^2V^*$を用いて
$\Pi=\Pi(F_R,a)$と表せることは同値である。
この同型は任意の$R$の基底変換と可換であり、
$\Pi$は発散零である。
\end{lemma}
\begin{proof}
$F$の五つの偏微分は二次正則列だから、そのKoszul複体の
\[
 \wedge^3V\otimes R
 \xrightarrow{\iota_{dF_R}}
 S_{R,2}\otimes\wedge^2V
 \xrightarrow{\iota_{dF_R}}
 S_{R,4}\otimes V
\]
の該当次数は、特殊ファイバーで左端の単射と中央での完全性を持つ。
有限自由加群の複体なので、分裂小行列式の単元性を
Artin基底へ持ち上げて同じ完全性を得る。
$\wedge^3V\simeq\wedge^2V^*$により表示と一意性が従う。
$d(dF_R\wedge a)=0$から発散零である。
\end{proof}

\begin{proposition}[JacobiイデアルとPl\"uckerイデアルの一致]
\label{prop:artin-pfaffian}
$\Pi=\Pi(F_R,a)$に対し、Jacobi方程式の係数イデアルは
\[
 (q_0(a),\ldots,q_4(a))
\]
と一致する。特に集合論的なrank分類だけでなく、
非被約な局所Artin基底上での方程式を同定する。
\end{proposition}
\begin{proof}
$J=\frac12[\Pi,\Pi]$と書くと、五変数の外積計算から
\[
 \iota_{\Pi\wedge\Pi}\Omega=2g\,dF_R,\qquad
 \iota_J\Omega=dg\wedge dF_R
\]
が得られる。これは多項式恒等式である。
符号は$ a=dx_0\wedge dx_1+dx_2\wedge dx_3$
の代入で直接確認できる。
線形写像
\[
 L_F:k^5\longrightarrow\Omega^2_{S,3},\qquad
 (c_r)\longmapsto d\bigl(\textstyle\sum_rc_rF_r\bigr)\wedge dF
\]
は単射である。核の元について、
$F_i$のKoszul完全性から
$d(\sum c_rF_r)$は$dF$の多項式倍となるが、
係数の次数は負なので零である。
Euler式と定数係数のKoszul syzygyの消滅により$c_r=0$。
従って$L_F$の非零$5$次小行列式は
$R$上でも単元であり、$L_{F_R}$は分裂単射となる。
上の$J$の表示の係数から$q_r(a)$を回収でき、
両方のイデアルが一致する。
\end{proof}

\begin{lemma}[同次量子化による中心三次商の構成]
\label{lem:artin-homogeneous-casimir-quantization}
$\Pi\in\mathfrak m\otimes S_2\otimes\wedge^2V$を
$[\Pi,\Pi]=0$を満たす二次bivectorとする。
Kontsevichの多重線形写像を一つ固定すると、
$R\otimes S$上に次数を保つ単位的な星積$B_\Pi$が定まり、
$B_\Pi$は一次生成の二次代数である。
この構成は任意のArtin係数環の準同型と可換である。

twisted Taylor mapの次数三部分
$D^0_\Pi:S_{R,3}\to(B_\Pi)_3$は同型であり、
任意の$F_R\in S_{R,3}$について
\[
 D^0_\Pi(F_R)\in Z(B_\Pi)
 \quad\Longleftrightarrow\quad [\Pi,F_R]=0.
\]
特に$F_R\bmod\mathfrak m=F$および$q_r(a)=0$なら、
$\Pi=\Pi(F_R,a)$、$w_R=D^0_\Pi(F_R)$から
次数別有限自由な中心三次商$A_R=B_\Pi/(w_R)$を得る。
\end{lemma}
\begin{proof}
Kontsevichの写像$U_n$をCattaneo--Felder--Tomassiniの規約で取り、
$\alpha=\Pi/2$として
\[
 f\star g=fg+\sum_{n\ge1}\frac1{n!}
 U_n(\alpha,\ldots,\alpha)(f,g),\qquad
 D^0_\Pi(f)=\sum_{n\ge0}\frac1{n!}
 U_{n+1}(f,\alpha,\ldots,\alpha)
\]
と定める。$\Pi$の係数は$\mathfrak m$に入るので和は有限である。
結合律と単位はformality恒等式から従い、一次交換子は$\Pi$である。
ここで用いる正確な外部入力はCattaneo--Felder--Tomassini
\cite[\S3, Theorem~3.1, Proposition~3.3]{CFT}であり、
その式はテンソルの係数と微分についての普遍的多項式である。
従って実数から複素数へ係数拡大し、冪零Artin係数を代入できる。
\S3で述べられているGL同変性もこの係数拡大後に成立する。

各グラフでは二次bivectorの頂点から二本の辺が出て、
各辺は一回の微分を与える。
従って星積は通常の多項式次数を保ち、
$D^0_\Pi$は三次式を三次式へ送る。
基礎加群は$R\otimes S$なので次数別有限自由であり、
一次生成性は各次数でのNakayamaによる。
一次生成空間からのテンソル代数の核を$K$とすると、
$K_n$は有限自由な直和因子で、その特殊化は多項式環の関係イデアルである。
$K_2$から生成されるイデアルは特殊化後に各$K_n$を生成するから、
$K_n$へのNakayamaにより$K=(K_2)$となる。
基底変換に関する自然性は、上の有限の普遍式から従う。

$D^0_\Pi$の次数三部分
$D^0_\Pi:S_{R,3}\to(B_\Pi)_3$は
還元が恒等の有限自由加群の同型である。
また係数が四次のベクトル場$Y$に対して
\[
 D^1_\Pi(Y):=\sum_{n\ge0}\frac1{n!}
 U_{n+1}(Y,\alpha,\ldots,\alpha)
\]
とおく。その評価写像
\[
 Y\longmapsto
 \bigl(D^1_\Pi(Y)(x_0),\ldots,D^1_\Pi(Y)(x_4)\bigr)
\]
は$S_{R,4}\otimes V\to(B_\Pi)_4^{\oplus5}$の写像であり、
同じく還元が恒等の有限自由加群の同型である。
Cattaneo--Felder--Tomassini
\cite[Corollary~5.2(iv)]{CFT}のtwisted chain identityにより
\[
 D^0_\Pi(F_R)\in Z(B_\Pi)
 \quad\Longleftrightarrow\quad [\Pi,F_R]=0
\]
となる。ここで中心の形成と非平坦基底変換の交換性は
仮定せず、有限自由行列の可逆性によって逆向きを証明する。
命題\ref{prop:artin-pfaffian}により、$q_r(a)=0$なら
$\Pi(F_R,a)$はPoissonであり、定義から$[\Pi(F_R,a),F_R]=0$である。
よって$w_R$は中心元となり、補題\ref{lem:artin-normal-central}から
その正則性と商の次数別平坦性を得る。
この補題は具体的な族の構成であり、
次数付きHochschild変形と同次Poisson変形の全gauge分類を主張していない。
\end{proof}

\subsection{Ore局所化、偏極、および斉次二次包絡}
次数別平坦な斉次変形$A_R$を
$S/(F)$へ特殊化するものとする。
一次元の持上げ$\widetilde x_i$を選び、
標準開集合$U_i=D_+(x_i)\cap X$上で
\[
 \mathscr A_R(U_i)=
 \bigl(A_R[\widetilde x_i^{-1}]\bigr)_0
\]
を取る。$s=\widetilde x_i$、$\mathfrak m^N=0$とすると、
$[A_R,A_R]\subset\mathfrak m A_R$より
$\operatorname{ad}_s(\mathfrak m^rA_R)\subset\mathfrak m^{r+1}A_R$、
従って$\operatorname{ad}_s^N=0$である。
任意の$n\ge1$と$a\in A_R$に対し
\[
 s^{n+N-1}a=
 \sum_{j=0}^{N-1}\binom{n+N-1}{j}
 \operatorname{ad}_s^j(a)s^{n+N-1-j}\in A_Rs^n.
\]
逆側にも同じ展開があり、$\{s^n\}$は左右Ore条件を満たす。
$s$の特殊化は整域$S/(F)$の非零因子であり、
$A_R$の$\mathfrak m$進濾過の商における乗法も単射なので、
$s$は左右正則である。
局所化は$R$加群として$A_R$の乗法写像によるfiltered colimitだから平坦であり、
その各次数成分も直和因子として平坦である。
重なりは二つの一次元をともに可逆化した次数零部分で
記述され、三重重なりのcocycleも局所化の普遍性に従う。
この操作で$X$上に平坦な結合$R$代数層
$\mathscr A_R$が得られる。

\begin{proposition}[偏極とsection algebraの回復]
\label{prop:artin-ore-section}
上のOre局所化層の次数$n$部分を貼り合わせた
$\mathscr E_n$は可逆両側$\mathscr A_R$加群であり、
\[
 \mathscr E_m\otimes_{\mathscr A_R}\mathscr E_n
 \simeq\mathscr E_{m+n}.
\]
$n\ge0$について
\[
 H^i(X,\mathscr E_n)=0\ (i>0),\qquad
 (A_R)_n\xrightarrow{\sim}H^0(X,\mathscr E_n)
\]
が成り立ち、積とも両立し任意のArtin基底変換に可換である。
\end{proposition}
\begin{proof}
局所化の一次元は可逆なので強い$\mathbb Z$次数付けと
両側可逆性が従う。アフィン交叉上では、$\mathscr E_n$の有限な
$\mathfrak m$進濾過の各商は$\OO_X(n)$の有限直和となり、非負次数に限らず
高次コホモロジーを持たない。従って有限標準被覆の\v{C}ech複体で
大域コホモロジーを計算できる。
$R$の小拡大を順に帰納すると、
特殊ファイバーの$H^i(X,\OO_X(n))=0$ ($i>0,n\ge0$)
から高次消滅と$H^0$の有限自由性が得られる。
特殊化で同型な$(A_R)_n\to H^0(X,\mathscr E_n)$は
Nakayamaで同型となる。強次数付けにより積との両立も従う。
\end{proof}

\begin{lemma}[可逆両側偏極の同型類の一意性]
\label{lem:artin-polarization}
$\mathscr A_R$上の二つの平坦な可逆両側加群
$\mathscr M,\mathscr N$が同じ線束$L$に特殊化し、
$H^1(X,\OO_X)=H^0(X,T_X)=0$なら、
$\mathscr M\simeq\mathscr N$である。
$H^0(X,\OO_X)=k$なら各々の自己同型は$R^\times$である。
\end{lemma}
\begin{proof}
小拡大$R'\twoheadrightarrow R$の核$I$で
$\mathfrak mI=0$とする。
特殊ファイバーの局所基底の持上げにより
$\mathscr M,\mathscr N$は右加群として局所自由階数一。
同型の右加群としての持上げの障害は
$I\otimes H^1(X,\OO_X)$に属し消える。
二つの左作用の差は導分に対応し
$I\otimes H^0(X,T_X)=0$なので持上げは両側同型。
自己同型群は
$H^0(X,\OO_X)=k$と小拡大帰納から$R^\times$となる。
\end{proof}

\subsection{構造代数層の変形関手と座標標識の商}
\begin{lemma}[相対三次Poisson変形の形式局所環]
\label{lem:artin-relative-poisson-sheaf}
$\mathcal D_{\mathrm{Pois},X}(R)$を、$\OO_X$への還元同型を持つ
平坦な可換$R$代数層と、還元が零である$R$線形Poisson括弧の
gauge同値類の集合とする。するとこの関手は
\begin{equation}\label{eq:artin-poisson-sheaf-ring}
 \mathcal R_X=
 \frac{k[[t_1,\ldots,t_{10},a_{01},\ldots,a_{34}]]}
 {(q_0(a),\ldots,q_4(a))}
\end{equation}
で表される。ここでは特殊ファイバーの有限自己同型群$\Aut(X)$を割らない。
\end{lemma}
\begin{proof}
まずPoisson括弧を忘れる。滑らかなアフィン開集合上の
冪零な平坦可換変形とその貼合せにより、基礎の代数層は
$X$の滑らかなArtin変形$X_R$の構造層となる。
小拡大の核を$I$とすると、線束$\OO_X(1)$の持上げの障害は
$I\otimes H^2(X,\OO_X)$に、持上げの差は
$I\otimes H^1(X,\OO_X)$に属する。
両者が零なので、$\OO_X(1)$は各$X_R$へ同型を除いて一意に持ち上がる。
三次超曲面の標準完全列から
\[
 H^1(X,\OO_X(n))=H^2(X,\OO_X(n))=0\quad(n\in\mathbb Z)
\]
である。特に$n\ge0$では高次コホモロジーがすべて消えるので、
小拡大帰納により持ち上げた線束の非負次数のsectionは自由で、
任意のArtin基底変更と可換である。
その一次sectionの基底を選ぶと$X_R$は$\PP^4_R$に埋め込まれる。
実際、section環への写像は特殊ファイバーで全射なので
各次数でNakayamaにより全射であり、三次核は階数一である。
その生成元$F_R$の乗法は特殊ファイバーで単射だから分裂単射であり、
同じHilbert関数とNakayamaにより全核は$(F_R)$となる。
従って$X_R=V_+(F_R)$である。

相対随伴公式により
$\wedge^2T_{X_R/R}\simeq\Omega^1_{X_R/R}(2)$である。
相対余法完全列と制限完全列
\[
\begin{split}
0&\longrightarrow\OO_{X_R}(-1)
 \longrightarrow\Omega^1_{\PP^4_R/R}(2)|_{X_R}
 \longrightarrow\Omega^1_{X_R/R}(2)\longrightarrow0,\\
0&\longrightarrow\Omega^1_{\PP^4_R/R}(-1)
 \longrightarrow\Omega^1_{\PP^4_R/R}(2)
 \longrightarrow\Omega^1_{\PP^4_R/R}(2)|_{X_R}\longrightarrow0
\end{split}
\]
を用いる。左端二層の$H^0,H^1$は、超曲面完全列とEuler完全列から
零である。さらにEuler完全列の次数二部分から
\begin{equation}\label{eq:artin-relative-bivectors}
 H^0(X_R,\wedge^2T_{X_R/R})
 \simeq H^0(\PP^4_R,\Omega^1_{\PP^4_R/R}(2))
 \simeq R\otimes_k\wedge^2V^*
\end{equation}
を得る。これらは有限自由加群間の分裂完全列による計算なので、
非平坦なArtin基底変更とも可換である。
$F_R$の尺度と体積形式を固定した上で、この対応は
定数二形式$a$を$\iota_Ea$に送り、対応するbivectorは
$\iota_{\Pi}\Omega=dF_R\wedge a$の射影制限である。
これは偏微分$\partial_iF_R$が単元となる滑らかなchartで
随伴公式の縮約を計算して確認できる。

ここで$J=\frac12[\Pi,\Pi]$、$g=\sum_rq_r(a)\partial_rF_R$とする。
命題\ref{prop:artin-pfaffian}の恒等式をEuler場で縮約すると
\begin{equation}\label{eq:artin-intrinsic-jacobi}
 \iota_E\iota_J\Omega
   =2g\,dF_R-3F_R\,dg
\end{equation}
となる。$F_R=0$に制限しEuler方向を除くと、
内在的Jacobi障害は随伴同型
$\wedge^3T_{X_R/R}\simeq\OO_{X_R}(2)$の下で
$2g$に対応する（全体の符号は縮約の規約による）。
五つの二次式$\partial_rF_R$は$R\otimes S_2$内で
階数五の直和因子を生成する。
実際、特殊ファイバーでの定数係数の関係は
滑らかな$F$の偏微分正則列のKoszul完全性により零であり、
非零最大minorを$R$へ持ち上げればよい。
また$H^0(X_R,\OO_{X_R}(2))=R\otimes S_2$である。
従ってJacobi方程式の係数イデアルは、
集合論的にだけでなく$(q_0,\ldots,q_4)$と一致する。

座標を残すと、$[F_R]$の34係数と$a$の10係数に
この五つのPl\"ucker方程式を課した関手を得る。
滑らかな三次三様体は$H^0(X,T_X)=0$を満たすので、
特殊ファイバーで恒等な射影座標変更の作用は自由であり、
その24次元軌道を分離する形式sliceが存在する。
この作用は分解可能二形式の錐を保つので、$F_R$をsliceへ移す
一意の形式座標変更を$a$にも作用させれば、
残る方程式は同じPl\"ucker方程式である。
これが\eqref{eq:artin-poisson-sheaf-ring}を与える。
特殊ファイバーを恒等に保つ自己同型は、最初の非零な
$\mathfrak m$進差分が$H^0(X,T_X)$に属するため存在しない。
従ってここで必要なのは集合値のprorepresenting環である。
\end{proof}

$\mathcal D_{\mathrm{cp},X}(R)$を、次数別有限自由な二次代数$B_R$と、
その三次部分の階数一の直和因子で中心に含まれる$L_R$の対の同値類の集合とする。
特殊ファイバーの$(B_R,L_R)$と$(S,kF)$の同一視をデータに含め、
同値はこの同一視を保つ次数付き同型とする。
すなわち、一次生成空間の座標を選んだ表示から、
特殊ファイバーで恒等な射影座標変更を割った関手である。
次数スカラー倍による自己同型をrigidifyしても同値類の集合は変わらない。
$\mathcal D_{\mathrm{Ass},X}$は$\OO_X$への還元同型を保つ
平坦結合代数層のArtin変形のgauge同値類の関手とする。

\begin{lemma}[構成した中心商族と代数層変形関手]
\label{lem:artin-constructed-family-comparison}
補題\ref{lem:artin-relative-poisson-sheaf}の三次式の形式slice
$F_{\boldsymbol t}$を固定する。
$\mathcal R_X$のArtin点$(\boldsymbol t,a)$に対し、
補題\ref{lem:artin-homogeneous-casimir-quantization}の
$\Pi(F_{\boldsymbol t},a)$の量子化と中心三次商を取り、
さらにOre局所化を取ることで得る自然変換
\[
 \Psi:\operatorname{Hom}_{\mathrm{loc}}(\mathcal R_X,-)
       \longrightarrow\mathcal D_{\mathrm{Ass},X}
\]
は自然同型である。
\end{lemma}
\begin{proof}
量子化の有限な普遍式、中心元の構成、およびOre局所化の普遍性により、
$\Psi$は任意のArtin基底変換に自然な実際の写像である。
Yekutieli\cite[Corollary~11.2]{YekutieliTwisted}は、滑らかな$X$と
$H^1(X,\OO_X)=H^2(X,\OO_X)=0$の下で、
可換積も動かすPoisson代数層変形と結合代数層変形の
gauge同値類の自然な全単射を与える。
本稿では$k=\C\supset\mathbb R$であり、局所Artin代数は
同定理のパラメータ代数に含まれる。
通常の変形とtwisted変形の同値類の一致は
同論文Corollary~7.5により保証されている。
補題\ref{lem:artin-relative-poisson-sheaf}により、
標的$\mathcal D_{\mathrm{Ass},X}$も完備局所Noetherian環
$\mathcal R_X$で表される。

接写像の単射性を直接示す。
$R=k[\varepsilon]/(\varepsilon^2)$上の点$(\dot{\boldsymbol t},a)$が
$\Psi$によって自明な変形へ移るとする。
一次交換子はgauge不変であり、構成した星積から得る交換子の射影制限は
\[
 \wedge^2V^*\longrightarrow H^0(X,\wedge^2T_X),
 \qquad a\longmapsto\iota_Ea
\]
という\eqref{eq:artin-relative-bivectors}の同型である。
従って$a=0$である。
このとき量子化は通常の積であり、中心三次商とその局所化は
可換な三次変形$F_{\boldsymbol t}=0$そのものである。
これが結合代数層として自明なら、可換代数層としても自明だから、
三次式のsliceの接方向$\dot{\boldsymbol t}$は零である。
よって$d\Psi$は単射であり、両側の接空間は20次元なので同型である。
この議論は、Yekutieliの量子化が選んだ幾何的座標を保つことを仮定しない。

標的の表現環との同型を一つ選ぶと、$\Psi$は$\mathcal R_X$の
連続自己準同型を与え、余接空間上で同型となる。
随伴次数環は次数一で生成されるから、各次数成分への写像は全射である。
それらは同じ有限次元空間の自己写像なので同型となり、
完備性と分離性から元の環準同型も同型である。
従って実際の$\Psi$が全Artin基底で自然同型となる。
\end{proof}

\begin{theorem}[斉次環の全Artin変形の二次cover]
\label{thm:artin-quadratic-cover}
Ore局所化による自然変換
\[
 \Theta:\mathcal D_{\mathrm{cp},X}
       \longrightarrow\mathcal D_{\mathrm{Ass},X}
\]
は自然同型となる。
従って任意の次数別平坦な斉次変形$A_R$は
\[
 V_R=(A_R)_1,\quad
 K_{2,R}=\ker(V_R^{\otimes2}\to(A_R)_2),\quad
 B_R^{\mathrm{quad}}=T_R(V_R)/(K_{2,R})
\]
を用いて一意の中心三次対から回復される。
特に$B_R^{\mathrm{quad}}$は次数別$R$平坦で、
\[
 L_R=\ker\bigl((B_R^{\mathrm{quad}})_3\to(A_R)_3\bigr)
 \subset Z(B_R^{\mathrm{quad}})_3
\]
は階数一の直和因子であり、
$A_R\simeq B_R^{\mathrm{quad}}/(L_R)$。
この構成はArtin基底変換と可換である。
\end{theorem}
\begin{proof}
$\Theta$は上のOre局所化の構成により定義され、
特殊ファイバーを保つ次数付き同型は局所化層のgauge同型を誘導する。
補題\ref{lem:artin-constructed-family-comparison}の$\Psi$は、
具体的な中心三次対の族を経由するので、$\Theta$は全射である。

単射性を示すため、二つの中心対$(B_R,L_R)$、$(B'_R,L'_R)$の
商$A_R$、$A'_R$の局所化層がgauge同型であるとする。
この同型によって同じ代数層上で考えると、
両者の局所化次数一部分は$\OO_X(1)$の可逆両側持上げであり、
補題\ref{lem:artin-polarization}によって両側同型となる。
特殊ファイバーの同型はスカラーを調整して恒等に選べる。
この同型のテンソル冪と命題\ref{prop:artin-ore-section}により、
積と特殊ファイバーの標識を保つ次数付き同型$A_R\simeq A'_R$を得る。
中心三次商は二次関係を変えないので、$B_R$と$B'_R$はそれぞれの
商の二次包絡である。従って上の同型は$B_R\simeq B'_R$へ一意に持ち上がり、
三次核$L_R$を$L'_R$へ送る。これで$\Theta$の単射性を得る。

任意の斉次変形$A_R$についても、そのOre局所化層に
$\Theta$の逆像となる中心対$(B_R,L_R)$を取る。
局所化次数一の両側持上げの一意性とsection回復により、
$A_R$はこの中心対の商と次数付き同型になる。
この比較を行ってから二次包絡を取れば
$B_R^{\mathrm{quad}}\simeq B_R$を得るので、その平坦性は未知のまま仮定していない。
三次核は$L_R$であり、主張する復元が従う。
基底変換は、有限自由な次数一・二の加群間の分裂核と
局所化の普遍性により可換となる。

Ore局所化の一般的な参照先はYekutieli\cite{YekutieliAffine}である。
De~Deken--Lowen\cite[Theorem~4.5]{DeDekenLowen}は
Abel変形とsection $\mathbb Z$-代数変形を比較するが、
同定理だけから次数付き構造や1周期性を推論してはいない。
全Abel変形への移行は定理\ref{thm:abel-cubic-flat-envelope}で、
各線束対象の持上げの一意性を用いて別に行う。
\end{proof}

\begin{theorem}[中心三次対の全Artin形式局所環]
\label{thm:artin-central-pair}
中心三次対の古典的変形関手$\mathcal D_{\mathrm{cp},X}$は
\[
 \operatorname{Spf}\mathcal R_X
 =\operatorname{Spf}\frac{
 k[[t_1,\ldots,t_{10},a_{01},\ldots,a_{34}]]
 }{(q_0(a),\ldots,q_4(a))}
\]
で表され、形式局所次元は17、接空間次元は20である。
一次生成空間の射影座標標識を残す関手は、非標準的に
\[
 \widehat{\mathcal I}_{F}\simeq
 \operatorname{Spf}\frac{
  k[[s_1,\ldots,s_{34},a_{01},\ldots,a_{34}]]
 }{(q_0(a),\ldots,q_4(a))}
\]
で表され、形式局所次元は41、接空間次元は44である。
有限群$\Aut(X)$はこの記述では割らない。
\end{theorem}
\begin{proof}
最初の主張は定理\ref{thm:artin-quadratic-cover}、
YekutieliのCorollary~11.2、および
補題\ref{lem:artin-relative-poisson-sheaf}から従う。
これは中心対の全ての変形を同次formalityで分類することを
前提としない証明である。

射影座標標識を忘れるときの群を
\[
 G_1(R)=\ker\bigl(\PGL_5(R)\longrightarrow\PGL_5(k)\bigr)
\]
とする。標識は自由な一次生成空間の射影基底だから、
その選択に$G_1(R)$が推移的に作用する。
系\ref{cor:artin-scalar}により中心対の無限小自己同型は
次数スカラー倍だけであり、射影基底を変えない。
従って標識を残す関手は$\mathcal D_{\mathrm{cp},X}$上の
$G_1$-torsorである。
補題\ref{lem:artin-constructed-family-comparison}で構成した族の
固定した生成空間$R\otimes V^*$がこのtorsorのsectionを与える。
$G_1$は24次元の滑らかな形式群なので、
標識を残す関手の表現環は$\mathcal R_X[[u_1,\ldots,u_{24}]]$である。
$t_1,\ldots,t_{10},u_1,\ldots,u_{24}$を$s_1,\ldots,s_{34}$と書き換えればよい。

Pl\"ucker環は$\mathrm{Gr}(2,5)$のアフィン錐の原点完備化であり、
その次元は7、接空間次元は10である。
従って$10+7=17$、$10+10=20$、
座標標識を残すと$34+7=41$、$34+10=44$を得る。
\end{proof}

\begin{remark}[一次方向はすべて二次へ延長するわけではない]
\label{rem:artin-obstruction}
$a=\varepsilon(dx_0\wedge dx_1+dx_2\wedge dx_3)$を
$k[\varepsilon]/(\varepsilon^2)$上で考えると
$q_r(a)=0$だが、
$k[\varepsilon]/(\varepsilon^3)$上では$q_4(a)=\varepsilon^2$。
二次の補正$b$を$a+\varepsilon^2b$に加えてもこの項は消えない。
従って接空間20次元と局所次元17の差は実際の二次障害を表し、
全Artin比較を無障害性と混同してはいけない。
\end{remark}

\begin{remark}[明確に主張しない事柄]
\label{rem:artin-limits}
この全Artin同型は偏極付き\emph{古典的}変形関手の同定であり、
derived moduli stack、高次自己同型を保つgroupoid、
構造対象を忘れたKuznetsov成分の変形、
またはそこで構成される可換復元との一致を与えない。
無限小の中心対を得ることと、有限非可換パラメータで
滑らかなtailsを持つことは別問題である。
\end{remark}

\subsection{固定ambientからの忘却の四次元ファイバー}
\begin{proposition}[中心三次直線の射影束]
\label{prop:artin-fixed-ambient-fiber}
ある良い開集合$U$の各点$B$が二次Koszul AS正則CY五次元整域であり、
$\dim Z(B)_1=2$、$\dim Z(B)_3=5$かつ
中心の形成が基底変更と可換であると仮定する。
$E_1=Z(\mathscr B)_1$、$E_3=Z(\mathscr B)_3$とおけば
\[
 0\longrightarrow\Sym^3E_1\longrightarrow E_3
 \longrightarrow L\longrightarrow0,\qquad\rank L=1.
\]
中心三次直線のパラメータ空間は
$\PP_U(E_3)\to U$で各fiberは$\PP^4$。
$\PP(E_3)\setminus\PP(\Sym^3E_1)$は
$\Sym^3E_1\otimes L^{-1}$を作用群とするtorsorである。
同型類を取る場合、$\Gamma_B=\Aut_{\gr}(B)/k^\times$
が有限なら一般fiberの粗い双有理型は
$\PP Z(B)_3/\Gamma_B$となり、その次元は4である。
\end{proposition}
\begin{proof}
中心三次部のベクトル束の射影化で前半を得る。
$B$が整域なので中心三次の非零元は正則、
次数ごとの乗法の核は零で、商の各次数が局所自由となる。
$E_3\to L$に非零に写る直線はこの短完全列のsplittingに対応し、
splittingの差は$\Hom(L,\Sym^3E_1)$に属する。
同型類については、二次包絡と三次核の復元から
$B/(W)\simeq_{\gr}B/(W')$と
$[W'],[W]$が$\Gamma_B$軌道に属することが同値である。
Rosenlichtの有理商を$U$とincidenceへ適用し、
周囲$B$の軌道の上の空間が
$\PGL(V)\times_{\Gamma_B}\PP Z(B)_3$となることから
不変関数体が$k(\PP Z(B)_3)^{\Gamma_B}$に一致する。
有限群商は次元を保つ。
\end{proof}

\begin{remark}[四係数は標識を忘れても消えない]
特に固定した周囲に対し$Z(B)_3=kG\oplus\Sym^3Z(B)_1$なら、
$G+f(u,v)$の四係数が相対アフィンfiberを表す。
有限群$\Gamma_B$による同一視は生じ得るが、
四係数全体を座標変換で消せるわけではない。
原稿中の明示剛体周囲$B_*$についての
$\Aut_{\gr}(B_*)=k^\times$と正則開集合の主張は、
付属の有限整数証明書が現行Git treeに存在しないため、
本稿では独立に再検証した定理として採用しない。
\end{remark}
